\scanpage{23}

Mais alors ce système n'est autre que
\[
\overline S'=\overline S'(N',e',(y'_\varepsilon)_{0\leq\varepsilon\leq e'}).
\]
Notons que $p^e\geq N+1=(N'+1)p$ implique $p^{e-1}\geq N'+1$ i.e. $p^{e'}\geq N'+1$. Donc on est dans les $=$s conditions que précédemment, avec cette fois-ci $N'<N$. Par hyp. de récurrence on a donc $y'_0=0$, donc $y_0=0$, cqfd. Ouf !

\uline{Théorème 1} Soient $E$ un \unsure{groupe} abélien, $N\in\mathbb N$, $(x_i)_{i\in\mathbb N}$ une famille d'éléments de $E$, telle qu'on ait
\figure{\begin{tikzpicture}\draw (0,0)--(0,-1.4);\end{tikzpicture}}
\[
\binom ji(x_j-x_i)=0\qquad\text{pour }j>i\geq N
\]
Alors on a $x_i=x_N$ $\forall i\geq N$.

\uline{Dém} Quitte à localiser en tous les nbs premiers $p$, OPS $E$ un module sur \inserted{un} $\mathbb Z_p$. Il suffit de prouver $x_{N+1}=x_N$ (car par récurrence on tire de m $x_{N+1}=x_{N+2}$ etc). Choisissons $e$ \inserted{$\in\mathbb N$} tel que $p^e\geq N+1$, n'utilisons que les $x_i$ pour $i$ de la forme $N$ ou $N+p^\varepsilon$, $0\leq\varepsilon\leq e$. Le lemme 2 nous montre que $x_{N+1}=x_N$, cqfd.

\uline{Cor 1} Soient $k$ un anneau commutatif, $S=k\{T\}$ l'alg. des polynômes à puiss. div., \deleted{M un $k$-module, $\unsure{S,S^{+(N)}}$} \deleted{\unsure{M}}\inserted{$S^{+(N)}$} l'idéal à p.div. \deleted{engendré par $\unsure{T^{(N)}}$} défini par
\[
S^{+(N)}=\sum_{i\geq N}k.T^{(i)}.
\]
et considérons, pour $\forall$ $k$-module $M$, le $S$-module gradué
\[
\deleted{\unsure{a}}\quad\omega_N(M)=M\otimes_S S^{+(N)}=\sum_{i\geq N}M.T^{(i)}
\]
Alors :

\mspage{22}
